\chapter{Background and related work}
\label{ch:background}
\section{Trajectories as geometric and temporal objects}
A sampled trajectory is an ordered sequence $X=((t_i,p_i))_{i=0}^{T-1}$ with $p_i\in\mathbb{R}^2$. The polyline through its points defines geometry. The assignment of time to those points defines motion. Even on a straight road these properties can be separated. A vehicle covering one hundred metres in ten seconds can travel at constant speed or spend half the window waiting. Its endpoints and total displacement are identical in both cases, but its sequence of increments differs. A trajectory synthesizer therefore needs both a spatial support and a distribution over movement on that support.

A directed lane graph supplies information that a background map cannot convey by itself. For example, a bridge may lie directly above a second street in plan view. The distance between projected centerlines can be zero even though a vehicle cannot turn between them. At a junction, the permitted successor also depends on direction and participant class. A small distance to the nearest road is useful evidence about position. A connected assignment of road states throughout the sequence establishes a stronger property. Both depend on the accuracy and version of the map.

Microscopic simulation offers a practical test bed. The MoST scenario represents Monaco and its surroundings in SUMO, including multimodal roads and mobility demand \cite{most}. Simulation exports can be sampled frequently and can carry participant types without manual annotation. They inherit the modelling choices of the simulator. A generated trace resembling the export can therefore reproduce a simulated mobility pattern even when its performance on traffic measured in the field is unknown.

\section{Diffusion models for trajectories}
A denoising diffusion model learns to recover structure from samples corrupted at known noise levels \cite{ddpm}. The representation of the clean sample is a design choice. When the representation is a sequence of positions, denoising acts on a coordinate plane. In the present work it acts on the temporal sequence of distances travelled along a sampled path. Either representation can produce coordinates as output. Their intermediate states and their routes to enforcing feasibility differ.

DiffTraj models spatial and temporal structure in GPS trajectory synthesis \cite{difftraj}. Its use of diffusion demonstrates the value of modelling a distribution of trajectories when multiple samples can be valid. A coordinate sample does not by itself name a permitted successor lane. Diff-RNTraj addresses road-constrained generation by representing positions through segment information and continuous moving ratios, and learning to denoise a vectorized representation \cite{diffrn}. This is particularly relevant prior work: the general concept of road-constrained diffusion generation predates this study.

A related design, RNTrajGen, separates road-level movement from the production of finer trajectories \cite{rntrajgen}. These studies make a useful comparator at the level of modelling assumptions. The temporal-resolution pilot here supplies a permitted planned route and compares learned motion with empirical transfer on identical paths. The later system samples a route from a learned policy over a permitted catalog. Its final GAF comparison holds that route policy common, preserving attribution of motion changes while allowing alternative routes across requests and draws. Published systems trained on different cities, sampling rates and definitions of trajectory quality provide context rather than numerical head-to-head results.

Scenario Diffusion generates agents and their trajectories jointly under map and scenario conditioning \cite{scenario}. Its output represents a different unit: a traffic scene. In our experiment the same starting position can generate several independent participant traces. Simultaneous occupancy and interaction are therefore outside the output distribution. This is consequential for interpretation. A bus trace can follow connected lanes while failing to be compatible with the timing of a particular traffic signal or the trajectory of another vehicle.

\section{Imaging a time series}
Wang and Oates encode time series using Gramian angular fields and Markov transition fields \cite{wang}. A pixel represents a relation between temporal indices. In a Gramian field the relation is computed from angular transforms of the underlying scalar values. In a Markov transition field it depends on the estimated transition probabilities between discretized states. Convolutional networks can process the resulting arrays as images. The construction gives spatially local filters access to structured patches of pairwise relations.

For a traffic trajectory the input has two spatial coordinates. The inherited implementation maps each point to a scalar using a Hilbert traversal of a spatial grid. The mapping changes the representation of a location while preserving the order in which observations were made. A CNN therefore sees temporal relations between locations collapsed to a scalar index. Nearby points along the Hilbert traversal occupy neighboring spatial cells, but some neighboring cells can be far apart in traversal order. The mapping is directional in this sense. It can also lose distinctions between slightly different positions in one cell.

Salvadori investigated such images for traffic participant recognition, including choices of encoding and classifier \cite{salvadori}. Ferreira and colleagues report a related classification study on simulated traffic \cite{ferreira}. The available historical repository provides an implementation of an image construction \cite{original}. We use that source to specify the transform. The historical publication, the university thesis, the inherited code and the weights trained for this experiment are distinct pieces of evidence. Differences in splits, sampling horizon and data preparation can change reported scores even for closely related representations.

\section{Evaluation under multiple objectives}
A generative model can satisfy one objective and miss another. A route decoder can enforce topology for any input sequence of speeds. A generator that copies one speed profile can pass the geometric test while providing little new training information. A very diverse generator can produce implausible behavior. We therefore evaluate geometric consistency, one-dimensional motion distributions, candidate diversity and training utility separately. The common route and candidate budget make the comparison interpretable along each axis.

For distributional evaluation we use the first Wasserstein distance. In one dimension this is the area between two empirical cumulative distribution functions. Its units are those of the measured quantity. A difference of twenty metres in the travel-distance comparison has a direct interpretation as a reduction in distribution discrepancy; it is not the error of a predicted endpoint for a particular participant. For downstream evaluation we use macro-F1 on held-out participants. The averaging treats a rare class as a separate classification task instead of letting the most frequent class dominate.

DDIM sampling provides a deterministic sequence of reverse updates conditional on the initial noise \cite{ddim}. Reusing paths, request order and random seeds reduces incidental changes in a comparison. Different temporal representations still have different noise dimensions, so a common seed is not an identical perturbation of the same vector. The temporal-resolution pilot reports two fixed model seeds on one new SUMO episode. The final GAF-conditioning comparison has four OD cohorts with paired windows and two model seeds. Its cohort intervals are descriptive; repeated windows, model seeds and candidate draws are not additional independent transport domains. Earlier resampling studies retain their own conditional scope.

\section{Learned representations and conditional synthesis}
Representation alignment and generative conditioning provide different interfaces between an encoder and a denoiser. REPA uses pretrained representations as an auxiliary training target for diffusion models~\cite{repa}. It motivates asking whether useful structure in a feature space can aid generative learning, but a class-trained traffic feature is not automatically an adequate target for within-class motion. Its performance must be established on generated trajectories, not inferred from image-domain findings.

Diffusion Autoencoders combine a learned encoder with a conditional diffusion decoder and fit a latent generative model for new sampling~\cite{diffae}. This supplies a relevant precedent for turning an observed representation into a sampleable condition. The present traffic adaptation uses two prescribed image transforms, partial CNN adaptation, a road/duration-conditioned code prior and a physical trajectory decoder. It does not claim the near-exact image reconstruction or inversion properties of the original architecture. A lossy image view can still be useful as a condition, but prior mismatch and an ignored condition are practical possibilities that require separate checks.

Route choice is a further component rather than an incidental coordinate error. Recursive models describe path preferences through local decisions and goal-dependent structure~\cite{routechoice}. Here, a masked next-edge policy learned from observed branching decisions scores a finite set of permitted routes. It is a behavioural cloning model with catalog renormalization, not an estimated recursive-logit equilibrium. This distinction makes both its learned preference and its coverage limitations explicit.

\section{What a trajectory generator learns}
A route model and a motion model face different ambiguities. TrajGAIL learns sequential vehicle decisions by adversarial imitation on a road network \cite{trajgail}. Such a policy can choose among candidate links, while the rate at which a vehicle travels along a link remains a separate quantity. Diff-RNTraj jointly represents segment identity and movement within segments \cite{diffrn}. The two-stage design of RNTrajGen likewise separates road sequences from fine movement \cite{rntrajgen}. These distinctions motivate the temporal pilot comparator here. Transferring an empirical speed sequence to the very same path gives the baseline every geometric advantage of the proposed model. Because both candidates complete the prescribed route, the relevant improvement is in motion distributions rather than in reaching a different travel distance.

Map-conditioned forecasting has developed representations for both geometry and interaction. VectorNet encodes polylines, while LaneGCN propagates information through lane connectivity~\cite{vectornet,lanegcn}. Trajectron++ models heterogeneous agents and feasible dynamics, and Scene Transformer jointly predicts agent futures~\cite{tractron,scenet}; Social GAN captures multimodal pedestrian interactions~\cite{socialgan}, whereas Motion Transformer combines intention localization with movement refinement~\cite{mtr}. These methods typically exploit observed histories and surrounding traffic to forecast a scene. Complete-journey synthesis instead begins from an OD request and participant information, requiring a distribution over route choice, duration and movement for the entire trip.

\section{Guiding a trajectory while it is being generated}
In sequential decision problems, Diffuser treats planning as denoising a trajectory, while Diffusion Policy generates action sequences conditioned on observations \cite{diffuser,policy}. MotionDiffuser is a closer comparison for map-based movement: it represents joint agent futures in a compressed space and guides samples using differentiable costs \cite{motiondiff}. The Hilbert transform used here quantizes coordinates and changes sequence length. A small change to an increment can leave every encoded pixel unchanged until a cell boundary is crossed. A coordinate-gradient perturbation therefore needs an additional surrogate or a different interface. The final learned-condition experiment preserves the discrete transform and differentiates CNN parameters during training, without requiring derivatives through Hilbert indexing.

Classifier guidance uses a classifier gradient during the reverse process, while classifier-free guidance combines conditional and unconditional model estimates \cite{cg,cfg}. Universal Guidance extends the range of external guidance functions used at inference \cite{univ}. These methods make different assumptions about the information available at noisy states. In particular, a class score trained on clean trajectories has a physical interpretation only after a provisional motion sequence has been decoded. Inverse-problem samplers face a related distinction. DDRM uses structure in linear measurement models, whereas diffusion posterior sampling approximates a measurement term for more general observations \cite{ddrm,dps}. Participant classifiers provide a discriminative preference over trajectories, motivating separate evaluation of the motion distributions induced by their scores.

Twisted Diffusion Sampling uses weighted particles to approximate conditional distributions~\cite{tds}, while Feynman--Kac steering formalizes how rewards, resampling times and potentials influence the resulting samples~\cite{fk}. Compared with terminal selection over a finite pool of completed trajectories, intermediate resampling reallocates computation toward promising partial solutions, potentially changing which trajectories survive. Its effectiveness depends on the reliability of noisy-state preferences: early resampling can discard viable motions, whereas late resampling leaves less opportunity to modify them. The main approach studied here instead supplies a learned global code throughout denoising, coupling the representation to motion generation through a consistent conditioning interface. Fixed-pool selection and latent conditioning are therefore evaluated as separate mechanisms.

Sampling budget itself can change conclusions. Improved DDPM studies learned reverse variances and sampling efficiency \cite{idh}; EDM separates choices of noise schedule, preconditioning and solver \cite{edm}; DPM-Solver constructs efficient numerical updates for diffusion sampling \cite{dpm}. These works motivate explicit cost accounting. Equalizing network calls controls one cost, but a denoiser call on a long physical-time sequence is more expensive than one on 257 bins. The temporal comparison matches optimization steps and candidates and reports measured runtime for each representation. The earlier GAF selection studies reuse the same eight candidates. The final direct-draw comparison uses four samples per request and records latent-prior and motion-denoiser work separately.

\section{Representation, fidelity and utility}
External representations can enter diffusion training through feature alignment. REPA aligns projected denoiser states with clean-image encoder features~\cite{repa}, while centered kernel alignment provides a way to compare relationships between learned representations~\cite{kornblith2019}. In trajectory generation, these ideas motivate complementary measurements of representation structure and delivered motion quality: feature similarity characterizes what the encoder retains, and physical evaluation measures how that information affects synthesized trajectories.

A diffusion model over time does not uniquely specify a trajectory task. TimeGrad uses an autoregressive probabilistic model for multivariate forecasting, while CSDI conditions on observed entries for imputation \cite{timegrad,csdi}. Those problems have observed context that a synthesis request may lack. DiffWave shows that an entire ordered waveform can itself be the generated object \cite{diffwave}. For road motion the temporal target used here is more structured: nonnegative increments give progressively ordered positions on a chosen path. A plausible marginal distribution of increments can still give an implausible sequence if their temporal ordering is wrong.

Recognition offers another perspective on that loss of information. Temporal convolutional models offer an alternative to recurrent sequence encoders \cite{tcn,gru}. InceptionTime applies an ensemble of deep convolutional networks to time-series classification; ROCKET and MiniRocket show that convolutional transforms with simple classifiers can also be competitive \cite{inception,rocket,minirocket}. Unsupervised multivariate sequence representations can learn reusable embeddings without training an image classifier \cite{mtst}. These alternatives motivate evaluating class-related structure through several representation families. Raw geometry and motion classifiers serve as information diagnostics in the data study. Independent physical measurements provide the primary generation outcomes, alongside class-recognition diagnostics.

Classifier two-sample tests ask whether a learner can distinguish populations \cite{c2st}. Distributional precision and recall instead separate sample quality from coverage \cite{prd,pr}. A selector can increase concentration in a recognizable part of the data distribution while decreasing the range of plausible behaviors available to the user. For that reason we compare variation across candidates, time near rest and steady movement in addition to speed and acceleration. Esteban and colleagues use synthetic time series to train predictors evaluated on real data \cite{tstr}. Downstream augmentation asks a further practical question whose result need not track either classifier confidence or a single distribution distance. The final experiment measures this use directly with fixed downstream training budgets and separate reference cohorts. Its observed-only comparison must account for the larger labelled corpus used upstream to train the generator; a GAF-specific claim additionally needs the baseline and raw-representation controls.

